
\section{Validation and numerical details}\label{app:numerics}
\subsection{Validations}
\begin{enumerate}
\item The response tables of Appendix~\ref{app:response} were computed two ways, by the Gauss sums of Eq.~\eqref{eq:gauss} and by direct codeword amplitudes, on thousands of frames at both distances, and the $d=3$ window was simulated as a 13-qubit state vector with the faults at their literal positions. All three agree on every case.
\item The two-fault $d=3$ pattern of Fig.~\ref{fig:d3-fault-expansion}
was checked with both doubled probability diagrams in bloqade-tsim~\cite{tsim} and
amplitude diagrams in pyzx\_param~\cite{pyzx_param}. The three-fault $d=5$ pattern of
Appendix~\ref{app:d5-fault-v2} was checked with pyzx\_param, using the spider cutting with like-term collection of Ref.~\cite{wan} followed by stabiliser decomposition. The results agree with $(A,B)=(1/4,1/4)$ at $d=3$ and $(1/16,1/16)$ at $d=5$.
These are floating-point scalar checks of fixed fault patterns, not
independent rational certificates for the full noisy circuits.
\item Every fixed-$p$ reconstruction passes the prime-product bound, an unused-prime check and reverse-order CRT, and the repeat at $p=10^{-3}$ with a new contraction order reproduces $A$, $B$ and $\PL$ exactly.
\item The degree-ten coefficients agree with the earlier calculation through degree seven, and multiplying the quotient back into $a(x)$ reproduces $e(x)$ through degree ten, Eq.~\eqref{eq:recursion}.
\end{enumerate}

\subsection{Published comparisons}
Clifft values are read from Fig.~6 of Ref.~\cite{clifft} and SOFT values from Fig.~2 of Ref.~\cite{soft}. At every shared $p$ the exact value lies within their sampling uncertainty. The $d=3$ series expansion also matches with state-vector data from~\cite{msc} in Fig.~\ref{fig:d3comparison}.

\subsection{Measured cost}
Table~\ref{tab:cost} lists the worker time per point, including final-response preparation, cache construction, growth contraction and assembly, on single-threaded C++ code executed on an Apple Macbook (M4)\footnote{In today's value: $<1000\text{USD}$.}. The $p=10^{-3}$ repeat, which reused the cached endpoint tables of Appendix~\ref{app:F-reuse}, took 1.710 worker core-hours and the complete degree-ten calculation 9.857 worker core-hours.

{\color{black}
\begin{table*}[b]\centering\small
\captionsetup{position=above,skip=4pt}
\caption{Prime counts are reconstruction plus held-out checks; bound bits refer to the sufficient full denominator.}
\label{tab:cost}
\begin{adjustbox}{max width=\linewidth}\begin{tabular}{rllll}
\toprule
$p$ & Bound bits & Endpoint primes & Full primes & Worker core-hours \\
\midrule
$0.0005$ & 21,137 & 124+2 & 347+1 & 11.090 \\
$0.00075$ & 20,477 & 118+2 & 336+1 & 10.434 \\
$0.001$ & 19,737 & \emph{reused} & 324+1 & 1.710 \\
$0.0015$ & 18,186 & 107+2 & 299+1 & 9.297 \\
$0.002$ & 18,080 & 101+2 & 297+1 & 8.642 \\
\bottomrule
\end{tabular}
\end{adjustbox}
\end{table*}}

\subsection{Exact coefficients}
Table~\ref{tab:d3coeff} and ~\ref{tab:d3raw} gives the $d=3$ series expanded coefficients, while Table~\ref{tab:d5raw} gives the $d=5$ numbers.
\begin{table}[H]\centering\small
\captionsetup{position=above,skip=4pt}
\caption{The $d=3$ coefficients in $P^{(3)}_\text{L}(x)=\sum_kL_kx^k$.}
\label{tab:d3coeff}
\begin{tabular}{r@{\qquad}l}
\toprule
$k$ & $L_k$ \\
\midrule
2 & $\frac{32}{75}$ \\[1.2ex]
3 & $\frac{2009039}{4500}$ \\[1.2ex]
4 & $\frac{1183791011}{54000}$ \\[1.2ex]
5 & $\frac{5576522343343}{6075000}$ \\[1.2ex]
6 & $\frac{7072572702107929}{243000000}$ \\[1.2ex]
7 & $\frac{3056416241583144971}{3645000000}$ \\[1.2ex]
8 & $\frac{14000422301261560415557}{656100000000}$ \\[1.2ex]
9 & $\frac{352460798986209480998423}{820125000000}$ \\[1.2ex]
10 & $\frac{2800398004859242449492508511}{590490000000000}$ \\[1.2ex]
\bottomrule
\end{tabular}

\end{table}

\begin{table*}[p]\centering\small
\captionsetup{position=above,skip=4pt}
\caption{The $d=3$ raw acceptance and accepted-error coefficients.}
\label{tab:d3raw}
\begin{minipage}[t]{.48\linewidth}\centering
\begin{tabular}{r@{\qquad}l}
\toprule
$k$ & $a_k$ \\
\midrule
0 & $1$ \\[1.2ex]
1 & $\frac{77}{3}$ \\[1.2ex]
2 & $\frac{109433}{100}$ \\[1.2ex]
3 & $\frac{203757811}{6750}$ \\[1.2ex]
4 & $\frac{7640477983}{8100}$ \\[1.2ex]
5 & $\frac{3345046003099}{121500}$ \\[1.2ex]
6 & $\frac{745483248873863}{911250}$ \\[1.2ex]
7 & $\frac{16316109725855291383}{683437500}$ \\[1.2ex]
8 & $\frac{2814706392354594541187}{4100625000}$ \\[1.2ex]
9 & $\frac{740868888180195867334063}{38443359375}$ \\[1.2ex]
10 & $\frac{2424452455579323062802516617}{4613203125000}$ \\[1.2ex]
\bottomrule
\end{tabular}

\end{minipage}\hfill
\begin{minipage}[t]{.48\linewidth}\centering
\begin{tabular}{r@{\qquad}l}
\toprule
$k$ & $e_k$ \\
\midrule
0 & $0$ \\[1.2ex]
1 & $0$ \\[1.2ex]
2 & $\frac{32}{75}$ \\[1.2ex]
3 & $\frac{2058319}{4500}$ \\[1.2ex]
4 & $\frac{3046313977}{90000}$ \\[1.2ex]
5 & $\frac{3010251653251}{1518750}$ \\[1.2ex]
6 & $\frac{8250007730158441}{91125000}$ \\[1.2ex]
7 & $\frac{223846253563431581}{60750000}$ \\[1.2ex]
8 & $\frac{695773372331253998653}{5125781250}$ \\[1.2ex]
9 & $\frac{473468329549333222967513}{102515625000}$ \\[1.2ex]
10 & $\frac{1347503441280814153010426947}{9226406250000}$ \\[1.2ex]
\bottomrule
\end{tabular}

\end{minipage}
\end{table*}

\begin{table*}[p]\centering\small
\captionsetup{position=above,skip=4pt}
\caption{Whole-circuit $d=5$ raw acceptance and accepted-error coefficients.}
\label{tab:d5raw}
\begin{minipage}[t]{.48\linewidth}\centering
\begin{tabular}{r@{\qquad}l}
\toprule
$k$ & $a_k$ \\
\midrule
0 & $1$ \\[1.2ex]
1 & $\frac{166}{3}$ \\[1.2ex]
2 & $\frac{874357}{180}$ \\[1.2ex]
3 & $\frac{803222513}{3375}$ \\[1.2ex]
4 & $\frac{461785319813}{36000}$ \\[1.2ex]
5 & $\frac{6813675679442473}{12150000}$ \\[1.2ex]
6 & $\frac{35673102687723752419}{1458000000}$ \\[1.2ex]
7 & $\frac{7127303166983391609601}{7290000000}$ \\[1.2ex]
8 & $\frac{25074527061623692681497017}{656100000000}$ \\[1.2ex]
9 & $\frac{7051054714711008821570505193}{4920750000000}$ \\[1.2ex]
10 & $\frac{7771271684969588874728818721617}{147622500000000}$ \\[1.2ex]
\bottomrule
\end{tabular}

\end{minipage}\hfill
\begin{minipage}[t]{.48\linewidth}\centering
\begin{tabular}{r@{\qquad}l}
\toprule
$k$ & $e_k$ \\
\midrule
0 & $0$ \\[1.2ex]
1 & $0$ \\[1.2ex]
2 & $0$ \\[1.2ex]
3 & $\frac{574}{375}$ \\[1.2ex]
4 & $\frac{1041564221}{810000}$ \\[1.2ex]
5 & $\frac{1012441365667}{1620000}$ \\[1.2ex]
6 & $\frac{6646695516492431}{97200000}$ \\[1.2ex]
7 & $\frac{267643169821800115741}{43740000000}$ \\[1.2ex]
8 & $\frac{70659679997414182779587}{164025000000}$ \\[1.2ex]
9 & $\frac{527877089790266980903761157}{19683000000000}$ \\[1.2ex]
10 & $\frac{98061225429543733697140528883}{65610000000000}$ \\[1.2ex]
\bottomrule
\end{tabular}

\end{minipage}
\end{table*}
